\documentclass[10pt,letterpaper]{article}
\usepackage{cornell-notes}

\title{Clause 29.05.06: Non Member Arithmetic}
\author{}
\date{}
\setDocTitle{Clause 29.05.06: Non Member Arithmetic}
\setDocSubtitle{C++ 2024}
\setDocOwner{C++ 2024 Cornell Notes}
\setDocDate{\today}
\setCornellCollection{C++ 2024 Cornell Notes}
\setCornellUnitType{Clause}
\setCornellUnitNumber{29.05.06}
\setCornellUnitTitle{Non Member Arithmetic}
\hypersetup{pdftitle={Clause 29.05.06: Non Member Arithmetic},pdfsubject={Cornell notes},pdfauthor={}}

\begin{document}
\maketitle
\begin{CornellOverviewBox}{Overview}
\textbf{Classification:} Standard library - Normative\\[2pt]
\textbf{Learning objective:} Understand the rules, terminology, and semantic consequences associated with Non-member arithmetic.
\end{CornellOverviewBox}\begin{CornellSummaryBox}{Summary}
{\sffamily\bfseries\color{Accent} Summary}\par\vspace{3pt}Section 29.5.6 focuses on Non-member arithmetic. Apply the specified interface together with its mandates, effects, result conditions, exception behavior, and complexity constraints. When applying it, preserve the distinction between mandatory requirements, recommendations, permissions, and behavior deliberately left undefined, unspecified, or implementation-defined.
\end{CornellSummaryBox}

\end{document}
